\item \textit{¿Por qué no es posible esta situación?} Se ha diseñado un nuevo motor diésel que incrementa la economía de combustible respecto a modelos previos. Los automóviles con este nuevo diseño son increíblemente solicitados. Dos características del motor son responsables de este aumentado ahorro de combustible: (1) el motor está totalmente fabricado de aluminio para reducir el peso del automóvil y (2) el escape del motor se emplea para precalentar el aire a $50^\circ\text{C}$ antes de que entre al cilindro para aumentar la temperatura final del gas comprimido. El motor tiene una \textit{razón de compresión}, es decir, el cociente del volumen inicial del aire con su volumen final después de la compresión, de 14.5. El proceso de compresión es adiabático y el aire se comporta como un gas ideal diatómico con $\gamma = 1.40$.
